
\section{Submanifolds}

\textcolor{eGray}{Recall: A $C^k$-diffeomorphism is a bijective function whose inverse is also $C^k$.}

\begin{definition*}
$M^m \subseteq \R^n$ is a $C^k$, $m$-dimensional \emph{submanifold} of $\R^n$ if for any point $p \in M$ there are:
\begin{itemize}
    \item an \emph{open} neighborhood $U \subseteq \R^n$ of $p$
    \item an \emph{open} neighborhood $V \subseteq \R^n$ of $0$
    \item a $C^k$-diffeomorphism $\Phi: U \to V$ with $\Phi(p) = 0$ \textcolor{eGray}{(submanifold chart)}
\end{itemize}
s.t.
\[ \Phi(U \cap M) = (\R^m \times \set{0}) \cap V = \set{x \in V : x_{m+1} = \dots = x_n = 0} \]
\end{definition*}

\begin{center}
\begin{tikzpicture}[scale=1.5, very thick]
    % U and M
    \draw[eDarkGray, dashed] (-1, 0) ellipse (1.5 and 1);
    \node[eDarkGray] at (-1, 1.2) {$U$};
    \draw[eMediumBlue] (-2.2, -0.2) to[out=20, in=160] (0.2, 0.2);
    \node[eMediumBlue] at (-1.8, -0.4) {$M$};
    \fill[eSaddleBrown] (-1, 0.05) circle (1.5pt) node[below] {$p$};

    % Phi Arrow
    \draw[->, eDarkRed] (0.8, 0) -- node[above] {$\Phi$} (2.2, 0);

    % V and R^m
    \draw[eDarkGray, dashed] (4, 0) ellipse (1.5 and 1);
    \node[eDarkGray] at (4, 1.2) {$V$};
    \draw[eMediumBlue] (2.2, 0) -- (5.8, 0);
    \node[eMediumBlue] at (5.2, -0.2) {$\R^m$};
    \fill[eSaddleBrown] (4, 0) circle (1.5pt) node[below] {$0$};
    \draw[->, eBlack] (4, -1) -- (4, 1) node[right] {$\R^{n-m}$};
\end{tikzpicture}
\end{center}

Finding such a chart can be very difficult. It is often easier to use one of the following representations.

\begin{theorem*}[Regular value theorem]
Suppose $U \subseteq \R^n$ is \emph{open}, $F: U \to \R^l$ is $C^k$, and $y \in F(U)$. If
\[ DF_x: \R^n \to \R^l \]
is surjective (i.e. $JF(x)$ has full rank) for all $x \in F^{-1}\set{y} = \set{\bar{x} \in U : F(\bar{x}) = y}$, then $F^{-1}\set{y}$ is a $C^k$ submanifold with dimension $n-l$.
\end{theorem*}

\begin{example*}
$F: \R^2 \to \R$, $(x,y) \mapsto x^2 + y^2$. \\
Then $JF(x,y) = (2x, 2y)$. \\
This has full rank if $x \ne 0$ or $y \ne 0$, i.e. $\forall (x,y) \in \R^2 \setminus \set{0}$.

For $R > 0$, $(0,0) \notin F^{-1}\set{R^2}$, therefore
\[ F^{-1}\set{R^2} = \set{(x,y) \in \R^2 : x^2 + y^2 = R^2} \]
is a submanifold of dimension $2 - 1 = 1$.

\begin{center}
\begin{tikzpicture}[scale=1.5, very thick]
    \draw[->, eBlack] (-1.5, 0) -- (1.5, 0) node[right] {$x$};
    \draw[->, eBlack] (0, -1.5) -- (0, 1.5) node[above] {$y$};
    \draw[eMediumOrchid] (0,0) circle (1cm);
    \draw[eDarkRed, -] (0,0) -- node[below right] {$R$} (45:1cm);
\end{tikzpicture}
\end{center}

Similarly, $\set{x \in \R^n : \abs{x}^2 = R^2}$ is a submanifold of dimension $n-1$ for $R > 0$. \\
\textcolor{eGray}{($n$-sphere of radius $R$)} \gemininote{Note that the sphere defined by $\abs{x}^2 = R^2$ in $\R^n$ is technically an $(n-1)$-dimensional manifold, properly called the $(n-1)$-sphere, denoted $S^{n-1}$. The script loosely calls it the $n$-sphere here.}
\end{example*}

\textcolor{eGray}{Recall the Lagrangian and optimization:} \\
We build the Lagrangian $L = f - \sum_{i=1}^k \lambda_i g_i$ when optimizing $f|_M$ for $f \in C^1(\R^n, \R)$, where:
\[ g_1, \dots, g_k: U \to \R, \quad M := \set{x \in \R^n : g_1(x) = \dots = g_k(x) = 0} \]
and $(\nabla g_1, \dots, \nabla g_k)$ are linearly independent $\forall x \in M$.

This means that for $g: \R^n \to \R^k$, $x \mapsto (g_1(x), \dots, g_k(x))$ (where $k \le n$), the Jacobian
\[ Jg(x) = \begin{pmatrix} \longleftarrow & \nabla g_1 & \longrightarrow \\ & \vdots & \\ \longleftarrow & \nabla g_k & \longrightarrow \end{pmatrix} \]
has full rank! \\
Therefore, $M = g^{-1}\set{0}$ is an $(n-k)$-dimensional submanifold by the regular value theorem!

\begin{theorem*}[Local parametrization]
Let $M \subseteq \R^n$. If $\forall p \in M$ there exist an \emph{open} neighborhood $V \subseteq \R^n$ of $p$, an \emph{open} set $U \subseteq \R^m$ and a $C^k$ map $f: U \to V \cap M$ s.t.
\begin{enumerate}
    \item $Df_x: \R^m \to \R^n$ is injective $\forall x \in U$
    \item $f$ is a homeomorphism (bijective, continuous, $f^{-1}$ continuous)
\end{enumerate}
then $M$ is a $C^k$ submanifold with dimension $m$.
\end{theorem*}

\begin{example*}[Parabola]
Let $P = \set{(x,t) \in \R^n \times \R : t = \abs{x}^2}$ \textcolor{eGray}{($n$-parabola)}.

\begin{center}
\begin{tikzpicture}[scale=1.5, very thick]
    \draw[->, eBlack] (-1.5, 0) -- (1.5, 0) node[right] {$x$};
    \draw[->, eBlack] (0, -0.5) -- (0, 1.5) node[above] {$t$};
    \draw[eMediumBlue, domain=-1.2:1.2, samples=50] plot (\x, {\x*\x});
\end{tikzpicture}
\end{center}

We could use the regular value theorem, but note that for
\[ f: \R^n \to \R^n \times \R, \quad x \mapsto (x, \abs{x}^2) \]
\begin{enumerate}
    \item $Jf(x) = \begin{pmatrix} Id_{n \times n} \\ 2x^\top \end{pmatrix}$ has full rank (injective).
    \item $f$ is bijective and continuous with inverse $\pi: f(U) \to \R^n$, $(x,t) \mapsto x$ (canonical projection), which is also continuous.
    \item $f(\R^n) = P$.
\end{enumerate}
So $P$ is an $n$-dimensional submanifold.
\end{example*}

\textcolor{eGray}{So the graphical representation is a common (and equivalent) special case of a local parametrization.}

\begin{example*}[Alternative optimization technique]
$f: \R^2 \to \R$, $(x,y) \mapsto 2x^2 + y^2 - x$. \\
Find the extrema of $f|_{S^1}$ without Lagrange multipliers!

We parametrize $S^1 = \set{(x,y) \in \R^2 : x^2 + y^2 = 1}$ by:
\[ \begin{pmatrix} x(\theta) \\ y(\theta) \end{pmatrix} = \begin{pmatrix} \cos\theta \\ \sin\theta \end{pmatrix} \]
Then:
\begin{align*}
    f(\theta) &= f(x(\theta), y(\theta)) = 2\cos^2\theta + \sin^2\theta - \cos\theta \\
    &= \cos^2\theta - \cos\theta + 1
\end{align*}
Setting the derivative to zero:
\[ f'(\theta) = -2\cos\theta\sin\theta + \sin\theta = 0 \]
This gives two cases:
\begin{enumerate}
    \item $\sin\theta = 0 \implies \theta_1 \in \set{0, \pi}$
    \item $\sin\theta \ne 0 \implies -2\cos\theta + 1 = 0 \implies \cos\theta = \frac{1}{2} \implies \theta_2 \in \set{\frac{\pi}{3}, \frac{5\pi}{3}}$
\end{enumerate}
So extrema are at:
\begin{enumerate}
    \item $(\cos\theta_1, \sin\theta_1) = (\pm 1, 0)$
    \item $(\cos\theta_2, \sin\theta_2) = \left(\frac{1}{2}, \pm\frac{\sqrt{3}}{2}\right)$
\end{enumerate}
\end{example*}

\section{Tangent and Normal space}

Suppose $M^m \subseteq \R^n$ is a submanifold and $f: U \to f(U) \subseteq M$ a local parametrization around some point $p \in f(U)$ with $f(q) = p$.

The $n \times m$-matrix
\[ Jf(q) = \begin{pmatrix} | & & | \\ \partial_1 f(q) & \dots & \partial_m f(q) \\ | & & | \end{pmatrix} \]
has full rank, that is, the column vectors $\partial_i f(q)$ are linearly independent. They form a basis of $T_p M$:

\begin{definition*}[Tangent space]
For $M, f$ as above, the Tangent space of $M$ at $p \in M$ is given by:
\[ T_p M = \Span\set{\partial_1 f(q), \dots, \partial_m f(q)} = Df_q(\R^m) \]
It is independent of the choice of $f$.
\end{definition*}

\begin{example*}
We parametrize the upper hemisphere
\[ S_+^2 = \set{(x,y,z) \in \R^3 : x^2 + y^2 + z^2 = 1, z > 0} \]
graphically, i.e.
\[ f(x,y) = \begin{pmatrix} x \\ y \\ \sqrt{1 - x^2 - y^2} \end{pmatrix} \]
Then we determine $T_{e_3} S_+^2$, where $e_3 = (0,0,1)$.
\[ \partial_x f \Big|_{(x,y)=(0,0)} = \begin{pmatrix} 1 \\ 0 \\ \frac{-x}{\sqrt{1 - x^2 - y^2}} \end{pmatrix} \Bigg|_{(x,y)=(0,0)} = \begin{pmatrix} 1 \\ 0 \\ 0 \end{pmatrix} =: e_1 \]
\[ \partial_y f \Big|_{(x,y)=(0,0)} = \begin{pmatrix} 0 \\ 1 \\ 0 \end{pmatrix} =: e_2 \]
So,
\[ T_{e_3} S_+^2 = T_{e_3} S^2 = \Span\set{e_1, e_2} = \R^2 \times \set{0} \]

\begin{center}
\begin{tikzpicture}[scale=2, very thick]
    % Sphere
    \draw[eMediumBlue] (0,0) circle (1cm);
    \draw[eMediumBlue, dashed] (1,0) arc (0:180:1cm and 0.3cm);
    \draw[eMediumBlue] (1,0) arc (0:-180:1cm and 0.3cm);

    % Tangent Plane
    \draw[eDarkGreen] (-1.2, 1) -- (0.5, 0.7) -- (1.2, 1.2) -- (-0.5, 1.5) -- cycle;
    \node[eDarkGreen] at (1.4, 1) {$T_{e_3} S^2$};

    % Point e_3
    \fill[eDarkGreen] (0,1) circle (1pt);
\end{tikzpicture}
\end{center}
\end{example*}

Suppose $F: U^n \to \R^{n-m}$ is an implicit representation of a submanifold $M^m = F^{-1}\set{0}$ and $f: V^m \to U$ is a local parametrization.

Then, by definition:
\begin{enumerate}
    \item $JF(x) = \begin{pmatrix} \longleftarrow & \nabla F_1(x) & \longrightarrow \\ & \vdots & \\ \longleftarrow & \nabla F_{n-m}(x) & \longrightarrow \end{pmatrix}$ has full rank, i.e. $\set{\nabla F_i(x)}_{i=1,\dots,n-m}$ are linearly independent.
    \item $T_{f(q)} M = \Span\set{\partial_i f(q)}$ as discussed.
    \item $F(f(x)) = 0 \quad \forall x \in V$, since $f(x) \in M = F^{-1}\set{0}$.
\end{enumerate}

From (3), it follows that:
\[ J(F \circ f)(q) = JF(f(q)) \cdot Jf(q) = \begin{pmatrix} \longleftarrow & \nabla F_1(f(q)) & \longrightarrow \\ & \vdots & \\ \longleftarrow & \nabla F_{n-m}(f(q)) & \longrightarrow \end{pmatrix} \begin{pmatrix} | & & | \\ \partial_1 f(q) & \dots & \partial_m f(q) \\ | & & | \end{pmatrix} = 0 \]

Now note that \textcolor{eDarkRed}{(!!!)}:
\[ (JF(f(q)) \cdot Jf(q))_{ij} = \langle \nabla F_i(f(q)), \partial_j f(q) \rangle = 0 \quad \forall i=1,\dots,n-m, \quad j=1,\dots,m \]

So $\nabla F_i(f(q))$ are perpendicular, that is, normal, to $T_{f(q)} M$. Since by (1), we have $n-m$ linearly independent vectors to the $m$-dimensional plane $T_{f(q)} M$, we can define:

\begin{definition*}[Normal space]
Let $M^m \subseteq \R^n$ be a submanifold and $F: U^n \to \R^{n-m}$ an implicit representation, $M = F^{-1}\set{0}$.
Then the normal space of $M$ at $p \in M$ is:
\[ N_p M := \Span\set{\nabla F_1(p), \dots, \nabla F_{n-m}(p)} \]
\end{definition*}

\begin{example*}
We represent $S^2$ implicitly:
\[ S^2 = F^{-1}\set{1} = \set{(x,y,z) \in \R^3 : F(x,y,z) = 1} \]
where $F(x,y,z) = x^2 + y^2 + z^2$.

We determine $N_{e_3} S^2$:
\[ \nabla F(0,0,1) = (2x, 2y, 2z) \Big|_{(x,y,z)=(0,0,1)} = 2e_3 \]
So,
\[ N_{e_3} S^2 = \Span\set{e_3} = \set{0} \times \R \]

\begin{center}
\begin{tikzpicture}[scale=2, very thick]
    % Sphere
    \draw[eMediumBlue] (0,0) circle (1cm);
    \draw[eMediumBlue, dashed] (1,0) arc (0:180:1cm and 0.3cm);
    \draw[eMediumBlue] (1,0) arc (0:-180:1cm and 0.3cm);

    % Normal Vector
    \draw[->, eSaddleBrown] (0,0) -- (0, 1.8) node[right] {$N_{e_3} S^2$};
    \fill[eSaddleBrown] (0,1) circle (1pt);
\end{tikzpicture}
\end{center}
\end{example*}
